\section*{Strict actions of the packed telescope on presheaves}
\paragraph{Status and base.}
Issue 105; checked against the source at commit
\texttt{265d2a80b3db05ef7d797f776f13224dd0050edd}.
This tranche concerns the right side of \texttt{eq:telescope} and its
specific binary native comparison, not the strict span-algebra theorem.
Write $L=\mathcal P_V$ and $T=\coprod_{n\ge0}L^n$ on the ordinary category
$[B^{\mathrm{op}},\mathbf{Cat}]$, with strict natural transformations.
The existing \texttt{PackedFree.monad} supplies the length-zero unit
and multiplication $\mu$. Here $L$ acts on actual Grothendieck arrows.
An $L$-algebra is a presheaf $Z$ and an arbitrary natural
$a:LZ\to Z$. This $a$ is supplied, not interpreted automatically as
a contextad's dependent multiplication.
\paragraph{Source-label crosswalk.}
All unqualified identifiers below lie in
\texttt{ContextualComposition.GrothendieckStep.AccumulationAlgebra}.
Under \texttt{eq:telescope}, \texttt{eval,empty,singleton} give
$E_a:TZ\to Z$ by the pre-existing layer recursion, with zero layer
identity and one layer $a$. \texttt{flatten,algebra} prove on actual
natural transformations that $T E_a;E_a=\mu_Z;E_a$, and install $E_a$
as a mathlib Eilenberg--Moore algebra. \texttt{ordered} restricts the
equation to the ordered cut (inner $m$, outer $n$): transport by the
actual iterate nesting equality followed by $m+n$ recursive steps.
The inputs of $(1,0)$ and $(0,1)$ have the same total but different
outer tags. No claim that arbitrary accumulation preserves these
tags at the output. \texttt{singleton\_join,recovered,restrict,extend\_restrict}
use length one of the genuine multiplication to recover an arbitrary
$T$-algebra on its very carrier (literal equality, using its two laws).
\texttt{map\_eval,strictMap,hom} characterize strict homomorphisms in
both directions. \texttt{extend,readOne,unitComponent,counitComponent,unitIso,counitIso,triangle,equivalence}
give the category equivalence with a specified inverse: restrict the
structure map along the length-one generator, without changing the
carrier map. Both comparison isomorphisms are identities on carriers.
No fullness of a functor forgetting cuts is being presumed.
\paragraph{Native interface.} \texttt{atFlatten} takes the packed law at
a base object $A$, as equality of functors including their maps.
\texttt{nativeFlatten} then pastes this equality with the actual
\texttt{NativePacked.actualSquare}, postcomposed with $E_a(A)$.
The domain is the $Z$-weighted fibre over $A$ of the native sequential
product of two all-length spans. One route uses its ordered chart
equivalence, applies $T E_a$, then $E_a$. The other embeds the two-block
row, passes into the source fibre, uses the actual weighted flattening
functor there, exits the fibre, evaluates the resulting single chain
into $TZ(A)$, and finally applies $E_a(A)$. The comparison is an
isomorphism of functors on native arrows, not an equality of output
observables. All binary lengths and the moving internal boundary are
in scope. No global modification in $A$ for the native comparison
is asserted; the packed algebra is natural in $A$.
\paragraph{Size and split law.} For the packed portion the base has
object/hom universes $b,w$ independent of the grade and payload
universes $u,\max(u,v)$. The native fibre comparison inherits
$b=u,w=\max(u,v)$. Both parts take the specified split endpoint data
with identity/composition laws for the endpoint reindexing comparisons.
No invertibility of a general contextad endpoint unit or multiplication
is asserted or used.
\paragraph{Obstruction to a larger conclusion.} The structure map here
is $LZ\to Z$ on a presheaf. A strict algebra for the telescope monad
on spans has a map $\operatorname{Tel}(V)\to V$, with a different
length-one unit. A proof for presheaf carriers cannot be inverted to
recover native span maps, unit, binary accumulation, endpoint squares,
or normal-cleavage/strict-map laws. To reach the last sentence following
\texttt{eq:telescope}, construct the native algebra interface and both
recovery maps under those laws, then compare its prescribed dependent
accumulation with the presheaf action. Nor is the pseudomonad of
\texttt{eq:grothmonad} the monad $T$: it uses an input contextad's $I$,
$\varepsilon$, tensor and $\delta$. No general pseudomonad, weak
strictification, or representable recovery is proved here. The strict
packed equation covers its full nested input; the native witness is
for two arbitrary-length blocks, not a native arbitrary row of blocks.
An auxiliary pentagon before multiplication of four independently
tagged blocks is neither needed nor claimed by this construction.
\paragraph{Discriminating checks.} In \texttt{AccumulationProbe}, a
single-object base has Boolean endomaps. The one-step algebra projects
the grade arrow, while empty evaluation retains the payload arrow.
Opposite constant endomaps yield different outputs; a retained constant
is noninvertible. The two-cut paths are evaluated on actual nested
arrows and opposite placements of an empty cut retain different tags.
In \texttt{MovingAccumulation}, based on the previously reviewed
moving cut fixture, an internal boundary moves $4\to6$ while the
terminal context moves $6\to10$. At the free algebra, accumulation on
the actual weighted native domain distinguishes parallel arrows with
payload words $[4,9]$ and $[9,4]$; naturality of the new comparison
is applied to both. Faithfulness in this last test is special to
the chosen free action, not an extra assumption on $a$.
